\section{Introduction}

In Bernoulli bond percolation on the nearest-neighbor cubic lattice
$\mathbb Z^3$, each edge is independently open with probability $p$.
In Bernoulli site percolation, the independent open bits instead belong to
the vertices, and an open path must have every vertex open, including its
endpoints. Write $\mathbb P_p$ for the law of the model under consideration
and define its critical parameter by
\[
 p_c=\inf\{p\in[0,1]:
       \mathbb P_p(0\text{ belongs to an infinite open cluster})>0\}.
\]
By translation invariance and a countable union over vertices, the same
threshold is obtained by asking for the existence of any infinite cluster.
The bond and site thresholds are defined separately throughout.

\begin{theorem}\label{thm:main}
For nearest-neighbor Bernoulli bond percolation on $\mathbb Z^3$, there is
almost surely no infinite open cluster at the bond critical parameter.
The same assertion holds for nearest-neighbor Bernoulli site percolation
at the site critical parameter.
\end{theorem}

Write $\theta(p)=\mathbb P_p(0\text{ belongs to an infinite open cluster})$.
The theorem gives continuity at the transition: $\theta(p)$ tends to
zero as $p\downarrow p_c$. To see this consequence directly, for $n\ge1$ let $E_n$ be
the event that the origin reaches the boundary of the box
$[-n,n]^3\cap\mathbb Z^3$ by an open path inside that box. Then
$\mathbb P_p(E_n)\downarrow\theta(p)$ as $n\to\infty$. At $p_c$, the
theorem makes this limit zero. Given $\varepsilon>0$, first choose $n$
with $\mathbb P_{p_c}(E_n)<\varepsilon$; since $E_n$ depends on finitely
many bits, the same strict inequality holds for $p$ sufficiently close
to $p_c$. Thus $\theta(p)\le\mathbb P_p(E_n)<\varepsilon$ there.
The proof below establishes critical extinction for both models without
using an external criticality theorem.

\paragraph{History and the three-dimensional obstacle.}
Broadbent and Hammersley introduced percolation to study transport through
a random medium~\cite{BroadbentHammersley1957}. In two dimensions, Harris's
critical nonpercolation theorem and Kesten's identification of the square
bond threshold established the bond result~\cite{Harris1960,Kesten1980}.
Russo proved critical nonpercolation for nearest-neighbor site percolation
on the square lattice~\cite[Theorem~2(b)]{Russo1981}. These planar results
use geometry that is unavailable for paths in three dimensions.

In sufficiently high dimensions, Hara and Slade established mean-field
critical behavior through the lace expansion~\cite{HaraSlade1990}.
Fitzner and van der Hofstad obtained the nearest-neighbor bond range
$d\ge11$~\cite{FitznerHofstad2017}. Hara and Slade also noted that their method applies to site
percolation~\cite[p.~335]{HaraSlade1990}. Heydenreich and Matzke later gave
a detailed site lace-expansion treatment, proving the triangle condition
in sufficiently high dimensions and its critical nonpercolation
consequence~\cite[Theorem~1.1 and Section~1.3]{HeydenreichMatzke2020}. The explicit bond dimension bound is
not a site bound. Benjamini and Schramm formulated the broader criticality
question for quasi-transitive graphs with $p_c<1$~\cite[Conjecture~4]{BenjaminiSchramm1996}.

Grimmett and Marstrand proved convergence of slab thresholds to the
full-lattice threshold and developed a dynamic renormalization argument
using open seeds and sequentially controlled exploration~\cite{GrimmettMarstrand1990}.
Duminil-Copin, Sidoravicius, and Tassion subsequently proved critical
nonpercolation for bond percolation on slabs~\cite{DuminilCopinSidoraviciusTassion2016}.
As their discussion explains, extinction at every fixed slab threshold,
together with convergence of those thresholds, does not by itself give
critical extinction in the full cubic lattice. The proof here obtains the
required connections between finite regions from a finite comparison,
and fixes all scales before decreasing the percolation parameter.

Leder's public account of a Lean formalization reports critical bond
nonpercolation on every nearest-neighbor lattice $\mathbb Z^d$,
$d\ge2$~\cite{Leder2026}. Its first-contact comparison
\cite[Section~3.4, (GEN)]{Leder2026} already implies the finite joint connection inequality for
bonds: take the increasing cluster functional to indicate meeting the
target set. The subsequent public Kozma--Nitzan exposition explicitly
records this consequence and reports stronger finite
comparisons~\cite[``Conjecture~1 from the Lean proof of Conjecture~3'']{KNPublic2026}.
The related public site exposition reports the
independent labelled-hyperedge comparison, a conditioned-covariance
hierarchy, and critical site nonpercolation on $\mathbb Z^d$ for every
$d\ge3$~\cite[``One inequality for both models'',
``Conditional covariance inequalities'', and ``From (1) to the lattice'']{SitePublic2026}.
These two public expositions describe work produced by AI systems
prompted by Ahmed Bou-Rabee. Our purpose here is a
self-contained cubic argument, including the finite comparison, explicit
relay geometry, and exploration estimates.

\paragraph{Finite comparison and the proof method.}
Our main finite ingredient is a connection inequality for independent,
possibly nonuniform hyperedges. For a node $o$ and nonempty node set $A$,
it asserts that
\begin{equation}\label{intro:comparison}
 \mathbb P(o\leftrightarrow A,\ o\leftrightarrow T)
 \ge \mathbb P(o\leftrightarrow A)
          \min_{a\in A}\mathbb P(a\leftrightarrow T).
\end{equation}
Here an open hyperedge joins all of its nodes. The inequality implies that
if every possible relay in $A$ reaches a target $T$ with probability at
least $1-b$, then the probability of reaching $A$ but missing $T$ is at
most $b$. This failure bound also holds for sites, and after any set of bits
has been fixed while the remaining bits retain their product law. Thus it
remains available during an exploration.

Taking $T$ to be a singleton in \eqref{intro:comparison} and enlarging
the event on the left gives the multiplicative gluing inequality of
Kozma and Nitzan~\cite[Conjecture~1]{KozmaNitzan2024}.
They also record that van den Berg and Engelenburg independently
considered this inequality. Section~4 of Kozma and Nitzan connects finite comparisons
to criticality through renormalization. The proof of
\eqref{intro:comparison} uses increasing functions of clusters, a
nonnegative cone of conditional connection columns, and alternating
conditional cluster resampling. Conditional resampling has precedents in
percolation correlation inequalities \cite[Section~2.1]{BHK2006}.
We prove the required comparison and every lattice relay estimate locally.
The general bond criticality theorem does not imply the site assertion:
independent site bits become independent hyperedges, and their endpoint
convention must be preserved. The finite hyperedge comparison
and its stability under fixed bits are useful independently of the lattice
application.

\paragraph{Proof overview.}
Assume that an infinite cluster exists at $p_c$. Large seed boxes then
connect with high probability to each of the 24 quarter-faces of a surrounding
cube. Only finitely many radii are needed, so these estimates persist at one
parameter $q<p_c$.

The finite comparison lets us extend a connection through a fresh region.
The seed and sequential-exploration architecture is related to the dynamic
renormalization method of Grimmett and Marstrand~\cite[Sections~3--4]{GrimmettMarstrand1990}.
Here many well-separated entrance points provide independent chances for an
open seed, while relay quality is defined using the interior bits alone. A
sequence of thirteen elementary geometric moves transfers a connection from
the inner box of one coarse cube to that of its neighbor. Finally, we explore
cubes indexed by $\mathbb Z^2$, retaining a high conditional connection
probability for each pending cube. Each processed cube has a uniformly small
conditional chance of failure. A planar boundary count then gives a positive
chance of infinitely many successful cubes, and hence an infinite open
cluster at $q<p_c$, a contradiction. Neither uniqueness of the infinite
cluster nor independence of the successful coarse cubes is assumed.

Section~\ref{sec:comparison} proves the finite comparison.
Section~\ref{sec:preliminaries} gives the planar count and fixes the scales.
Sections~\ref{sec:extension} and~\ref{sec:geometry} prove the extension
and neighboring-box estimates. Section~\ref{sec:exploration} constructs
the exploration and completes the proof.

\paragraph{Connection conventions.}
For a vertex $x$ and vertex set $T$, the notation $x\leftrightarrow T$
means that an open path in the specified graph joins $x$ to some vertex of
$T$. For two sets, at least one pair of their vertices must be so joined.
We permit paths of zero edges. Such a path always gives a self-connection
for bonds and hyperedges, but for sites its single vertex must be open.
All graphs and hypergraphs in the finite comparison are undirected.
